\documentclass[
	% -- opções da classe memoir --
	12pt,				% tamanho da fonte
	openright,			% capítulos começam em pág ímpar (insere página vazia caso preciso)
	a4paper,			% tamanho do papel. 
	% -- opções da classe abntex2 --
	%chapter=TITLE,		% títulos de capítulos convertidos em letras maiúsculas
	%section=TITLE,		% títulos de seções convertidos em letras maiúsculas
	%subsection=TITLE,	% títulos de subseções convertidos em letras maiúsculas
	%subsubsection=TITLE,% títulos de subsubseções convertidos em letras maiúsculas
	% -- opções do pacote babel --
	english,			% idioma adicional para hifenização
	french,				% idioma adicional para hifenização
	spanish,			% idioma adicional para hifenização
	brazil				% o último idioma é o principal do documento
	]{article}

% Packages:
% \usepackage[alf]{abntex2cite}		% Citações padrão ABNT
\usepackage{lmodern} 				% Usa a fonte Latin Modern	
\usepackage[T1]{fontenc} 			% Selecao de codigos de fonte.
\usepackage[utf8]{inputenc}			% Codificacao do documento (conversão automática dos acentos)
\usepackage{lastpage} 				% Usado pela Ficha catalográfica
\usepackage{indentfirst} 			% Indenta o primeiro parágrafo de cada seção.
\usepackage{color} 					% Controle das cores
\usepackage{graphicx} 				% Inclusão de gráficos
\usepackage{microtype} 				% para melhorias de justificação
\usepackage{lipsum} 				% para geração de dummy text
\usepackage{booktabs}
\usepackage[table,xcdraw]{xcolor}
\usepackage{colortbl}
\usepackage{float}

\usepackage{cite}
\usepackage{hyperref}
\usepackage[brazilian,hyperpageref]{backref}

\bibliographystyle{plain}

\setlength{\parindent}{1.3cm}
\setlength{\parskip}{0.2cm}

\title{Plano de pesquisa para o Trabalho de Conclusão de Curso\\\large Análise comparativa de Modelos de Predição para Cheias no Rio Guaíba}
\author{Aluno: Juliano Almeida Machado \\
Orientador: Prof. Dr. Sebastián Gonçalves}
\date{}

\begin{document}
\maketitle

\section{Introdução}
Desastres naturais provocados por eventos climáticos extremos causam impactos profundos e severos de diversas maneiras. 
No aspecto econômico, os recursos alocados para a recuperação de áreas afetadas enfraquecem a microeconomia local, 
tornando a região ainda mais vulnerável a eventos futuros. Estruturalmente, o tempo e o dinheiro investidos na 
construção e no aprimoramento de edificações, tanto privadas (como residências e negócios) quanto públicas (como 
infraestrutura urbana e áreas administrativas), podem ser comprometidos pela falta de preparação para lidar com tais 
ocorrências. Já o impacto social é, sem dúvida, o mais significativo. Este impacto, imensurável, é sentido por diversas 
gerações de pessoas afetadas por desastres, como evidenciado pelos efeitos do furacão Katrina\cite{reardon}\cite{neria}.

Devido às mudanças climáticas, a frequência e intensidade dos desastres naturais vêm aumentando gradualmente desde os 
anos 1950\cite{dunn}. Como resultado, muitas regiões estão sendo afetadas por eventos climáticos extremos de maneira 
cada vez mais rápida e despreparada. Entre esses desastres, as enchentes se destacam pela sua capacidade de causar 
destruição generalizada e afetar áreas densamente povoadas, especialmente em regiões urbanas com infraestrutura 
inadequada para gerenciar grandes volumes de água. No final de abril e início de maio de 2024, o estado do Rio Grande do 
Sul passou por um período de alta precipitação que resultou na cheia das suas principais bacias hidrográficas. Na Região 
Metropolitana de Porto Alegre, esse cenário culminou na maior cheia histórica do Rio Guaíba, que ultrapassou a marca de 
4.76 metros, superando a cheia de 1941, até então considerada a mais devastadora da cidade\cite{dep}. Esse desastre foi 
impulsionado por uma combinação de fatores meteorológicos, e, combinado com a maior vulnerabilidade de centros urbanos à 
enchentes\cite{yang}\cite{hammond}\cite{zhang}\cite{pereira}, levou à um evento de devastação generalizada, com milhares 
de pessoas deslocadas e centenas de vítimas fatais\cite{alcantara}. 

Durante a gestão da crise, o monitoramento e predição dos níveis de água do Rio Guaíba – o mais importante da região – 
foi fundamental para orientar os tomadores de decisão em suas ações de planejamento e para que à população pudesse 
reagir melhor à ameaça iminente. Essas previsões foram feitas principalmente através de modelos hidrológicos e 
hidrodinâmicos, que, embora precisos, são complexos, exigem grande quantidade de dados topográficos, meteorológicos e 
hidrológicos, requerem conhecimento especializado em hidrologia e demandam tempo de simulação prolongado. Esse último 
fator foi acentuado nos momentos mais delicados da crise, quando as previsões só podiam ser realizadas com uma 
frequência diária \cite{iph}, o que as tornou insuficientes para atender às demandas de previsões em tempo real ou de 
curto prazo \cite{atashi}.

Uma alternativa que visa superar essas limitações são os modelos baseados em dados. Os mesmos utilizam técnicas como 
aprendizado de máquina para identificar e explorar relações estatísticas entre os dados de entrada e saída. Com a 
prendizagem dos padrões presentes nos dados históricos, é possível mapear as dinâmicas subjacentes do sistema e 
utilizá-las para prever valores futuros. Assim, esses modelos são capazes de realizar previsões precisas com base em 
séries temporais, capturando tendências e comportamentos complexos que podem ser difíceis de modelar diretamente por 
meio de abordagens tradicionais. Modelos baseados em dados são mais rápidos e eficientes para previsões em tempo real, 
sendo uma escolha adequada para sistemas de alerta precoce e mitigação de desastres de cheias\cite{ta}.

\section{Objetivos}
Este projeto tem como objetivo aplicar diferentes modelos de \textit{Machine Learning} nas séries históricas do nível da 
água do Rio Guaíba, enriquecidas com dados externos relevantes do mesmo período, como o nível de rios afluentes, 
precipitação em áreas críticas, velocidade e direção dos ventos, temperatura, entre outros. A inclusão desses dados 
externos será avaliada para determinar sua relevância, visando otimizar o modelo ao utilizar o menor número possível de 
variáveis. Os modelos serão treinados com os mesmos dados e avaliados com diversas métricas de desempenho. Será então 
realizada uma análise do comportamento e desempenho de cada modelo, buscando identificar padrões que possam contribuir 
para melhorar as previsões.

\section{Metodologia}
Neste trabalho, serão utilizados dados do Sistema Nacional de Informações sobre Recursos Hídricos (SNIRH), que 
disponibiliza informações das principais bacias hidrográficas do Brasil\cite{snirh}. Esses dados serão a base para as 
séries históricas do nível de água, complementados por dados externos, como os do MERRA-2 (\textit{Modern-Era 
Retrospective analysis for Research and Applications, Version 2}), um projeto da NASA que integra observações de 
satélites e estações terrestres para gerar dados meteorológicos globais\cite{nasa}. Dados meteorológicos adicionais 
serão obtidos a partir de fontes públicas, como o Instituto Nacional de Meteorologia (INMET) e o Instituto Brasileiro de 
Geografia e Estatística (IBGE)\cite{inmet}\cite{ibge}.

Esses dados serão processados em diferentes modelos de Redes Neurais Recorrentes, como \textit{Long Short Term Memory} 
(LSTM)\cite{hochreiter}, além de modelos de Perceptron Multicamadas (MLP) especializados, como \textit{Neural 
Hierarchical Interpolation for Time Series Forecasting} (NHITS)\cite{challu}, \textit{Neural Basis Expansion Analysis 
for Interpretable Time Series Forecasting} (NBEATS)\cite{oreshkin} e \textit{Time Series Mixer} (TSMixer)\cite{chen}. 
Todos esses modelos serão implementados em conjunto utilizando a biblioteca \textit{NeuralForecast}, desenvolvida pela 
empresa Nixtla\cite{neuralforecast}, que permite a execução paralela dos modelos sobre os mesmos dados. A validação 
cruzada por k-folds e as previsões dentro da amostra serão realizadas automaticamente pela estrutura da biblioteca.

A otimização dos hiperparâmetros será conduzida utilizando a ferramenta \textit{Tune} da biblioteca Ray\cite{tune}. Além 
disso, modelos baseados em árvores de decisão, como \textit{eXtreme Gradient Boosting} (XGBoost)\cite{chen_xgboost}, 
serão implementados paralelamente aos modelos mencionados. A comparação entre todos os  modelos será feita com base em 
métricas de desempenho como o \textit{Root Mean Squared Error} (RMSE), \textit{Mean Absolute Error} (MAE) e o 
\textit{Nash-Sutcliffe Model Efficiency Coefficient} (NSE).
\bigskip

\section{Cronograma}

\begin{table}[H]
\begin{tabular}{llllll}
\hline
                        & OUT/24                                          & NOV/24                                          & DEZ/24                                          & JAN/25                                          & FEV/25                   \\ \hline
Revisão Bibliográfica & \cellcolor[HTML]{656565}{\color[HTML]{333333} } & \cellcolor[HTML]{656565}{\color[HTML]{333333} } & \cellcolor[HTML]{656565}{\color[HTML]{333333} } & \cellcolor[HTML]{656565}{\color[HTML]{333333} } & \cellcolor[HTML]{656565} \\ \hline
Modelagem matemática    &                                                 & \cellcolor[HTML]{656565}{\color[HTML]{333333} } & \cellcolor[HTML]{656565}                        &                                                 &                          \\ \hline
Análise dos resultados  &                                                 & \cellcolor[HTML]{656565}{\color[HTML]{333333} } & \cellcolor[HTML]{656565}{\color[HTML]{333333} } & \cellcolor[HTML]{656565}                        &                          \\ \hline
Elaboração do texto     &                                                 &                                                 & \cellcolor[HTML]{656565}{\color[HTML]{333333} } & \cellcolor[HTML]{656565}{\color[HTML]{333333} } & \cellcolor[HTML]{656565} \\ \hline
Ajustes                 &                                                 &                                                 &                                                 &                                                 & \cellcolor[HTML]{656565} \\ \hline
\end{tabular}
\caption{Tabela do cronograma}
\label{table:1}
\end{table}

\bibliography{references}

\end{document}